\section{Choosing the parameters}\label{sec:parameters}

The preceding splits require two properties of their finite schedules.
For each aligned stage, the number of ambiguous cut offsets must be
smaller than the number of copies of every cut kind. At the same time,
every prefix and suffix test must see its prescribed data, including
the bounded controls at all false candidate origins. We now choose
the schedules so that both properties hold.

The number of cut kinds and the positions of their copies play different
roles. Enlarging the stage-0 roster creates more main trials and tags,
but the clock obstruction bound grows only logarithmically with the
tag count. This lets us fix the number of copies first. Their integer
positions remain free until the clocks and marker windows have fixed
the required main-trial width. The order of choices below makes this
separation explicit.

Throughout this section \(q\) is the anchor of the true episode.
Evaluating \(A_j(t+\delta,q)\) keeps that anchor fixed; it never
starts a new episode-end search at \(t+\delta\). This convention is
needed both for the clock estimates and for finite final padding.

\begin{proposition}\label{prop:parameters}
For every fixed nonempty finite alphabet, finite monoid \(M\), and
finite representation space \(V\), all parameters in the two outer
splits and the five aligned stages can be chosen to satisfy their
stated hypotheses. More precisely, the stencils are nonempty, have
more than \(C_i\) copies of each required kind, and satisfy
\[
 \#\{(p,p')\in\mathcal P_i^2:p\ne p',\
       A_i^{\Gamma_i}(t+p-p',q)\}\le C_i
       \qquad(0\le i\le4)
\]
on every complete episode. Write \(N_i=|\mathcal P_i|\).
In order, these bounds concern first-clock or prediction failure,
inner-search instability, second-clock failure, a right-margin marker,
and end-search instability. The constants can be taken to be
\begin{align}
 C_0&=K+c_{\rm cl}\log(2m_1),\label{eq:parameter-c0}\\
 C_1&=10|\mathcal Z|N_0^4(2R_1+3),\label{eq:parameter-c1}\\
 C_2&=N_0^4N_1^4c_{\rm cl}\log(2m_2),\label{eq:parameter-c2}\\
 C_3&=10|\mathcal Z|(2R_2+3)\prod_{j<3}N_j^4,
       \label{eq:parameter-c3}\\
 C_4&=10(2w_E+3)\prod_{j<4}N_j^4.
       \label{eq:parameter-c4}
\end{align}
Here \(K=|M|^{|V|^2+1}\), \(m_1,m_2\) are the respective full tag
alphabet sizes, and \(c_{\rm cl}\) is the absolute constant of
Lemma~\ref{lem:clock}; logarithms are natural. Also,
\(\operatorname{diam}\mathcal P_i<w_E\) for \(i<4\).
All measurements, trials, and inspections have the chronological and
reading-side access required in Lemmas~\ref{lem:main-split},
\ref{lem:periodic-split}, and~\ref{lem:aligned-stage}. Every bounded
control, including every control at a false candidate origin, lies
inside the true first episode. The same choices satisfy the end
margins of Section~\ref{sec:geometry}.
\end{proposition}

\begin{proof}
Write \(\kappa_i\) for the number of cut kinds at stage \(i\), and
use \(\mu_i\) copies of each kind, so that
\(N_i=\kappa_i\mu_i\). We will choose the finite objects in the
following order and verify each estimate as its inputs become fixed.
\begin{enumerate}
\item Fix the spaces \(U_i\), full trees, aligned trial identities,
      and prediction assignments. These determine the numbers
      \(\kappa_i\) before any clock or position is chosen.
\item Choose the stage-0 copy count \(\mu_0\) and its abstract roster.
      This fixes the main-trial identities and predictions, the first
      tag count \(m_1\), and \(\mathcal Z\). Choose the first clock
      and \(R_1\), retaining its prescribed roster residues but
      leaving the integer stage-0 positions open.
\item Choose the stage-1 count and numerical stencil. Then choose
      the inner marker separation \(d_\circ\), block length
      \(K_\circ\), and inspection radius \(r_\circ\).
\item Fix the second tag alphabet, choose the stage-2 count, and
      construct the second clock and numerical stencil. Then fix
      the periodic-window width \(R_2\).
\item Choose the stage-3 count and numerical stencil. Then choose
      the inner search halfwidth \(H\) and a common main-trial
      template large enough for all the prescribed inspections.
\item Place the stage-0 stencil with its reserved residues and with
      differences separated by more than this template length.
      Place the main trials; this determines the outer cut sets.
\item Choose the end-stability radius \(w_E\), then the stage-4
      count and numerical stencil, and finally the end marker
      parameters \(d_\bullet,K_\bullet,r_\bullet\).
\item Translate the aligned arrangements into reverse stage order.
      Choose \(L_0\), then \(o_\bullet\), then \(B_{\rm tail}\)
      to put every bounded inspection inside the true episode.
\end{enumerate}

At steps 3--5 the stage-0 positions, and hence the numerical sets
\(\Gamma_1,\Gamma_2,\Gamma_3\), are not yet fixed. Their cardinalities are already
bounded in terms of the chosen roster sizes. The clock and marker
estimates are uniform in the translated origin and in the nominal
search centers, so these cardinality bounds suffice until the actual
positions are selected. The only copy count that enters its own
obstruction bound is \(\mu_0\); its logarithmic dependence will be
resolved explicitly below.

\subsection*{Finite kinds before numerical parameters}
Put \(m=|M|\). The sets
\[
 U_0=V\oplus\mathbf F_2,
 \qquad U_{i+1}=\mathbf F_2^{\,U_i\times M}\quad(0\le i\le4)
\]
are fixed now; in particular
\( |U_{i+1}|=2^{m|U_i|}\).
At stage \(i\), set
\[
 J_i^{\rm al}=|U_{i+1}|^2,
 \qquad K_i^{\rm al}=m^{J_i^{\rm al}+1},
 \qquad b_i=K_i^{\rm al}+1.
\]
Use the full rooted \(m\)-ary tree with \(b_i\) measurements along each
root-to-leaf path. It has
\[
 e_i=\sum_{j=0}^{b_i-1}m^j
 \quad\hbox{measurement nodes},\qquad
 \ell_i=m^{b_i}\quad\hbox{leaves}.
\]
These formulas also apply when \(m=1\). For integers \(a,r\ge0\),
write \((a)_r=a(a-1)\cdots(a-r+1)\) when \(r\le a\), and
\((a)_r=0\) when \(r>a\). Allocate one aligned trial for each ordered
list of \(K_i^{\rm al}\) distinct leaves, and one spare trial. There
are \(\tau_i=1+(\ell_i)_{K_i^{\rm al}}\) trials.
On a listed trial, assign the \(K_i^{\rm al}\) possible product tuples
bijectively to its listed leaves; assign arbitrary tuples to all
other leaves. The spare trial has arbitrary predictions. Thus every
history, including every infeasible history, has a prediction on every
trial, as required in Lemmas~\ref{lem:total-tables}
and~\ref{lem:bad-histories}.

There is one measurement kind for each node and payload in
\(U_{i+1}\), and one trial kind for each trial, block, and leaf. Hence
the number of kinds introduced above is
\begin{equation}\label{eq:parameter-kinds}
 \kappa_i=e_i|U_{i+1}|+
              \tau_iJ_i^{\rm al}\ell_i.
\end{equation}
None of these sets or counts involves a clock, a marker radius, a
main-trial identity, or a numerical position. In the terminal map
\([v,a]\mapsto[\Phi_i(a)(v),1_M]\), the argument set \(U_i\times M\)
and the output basis set are already fixed. A later, more complicated
formula for the total table \(\Phi_i\) therefore does not introduce
new kinds or enlarge the spaces. The formulas in
Lemma~\ref{lem:total-tables} evaluate every hypothetical product; they
do not add a prefix-feasibility field to a kind.

We fix an ordering of these kinds and their eventual copies compatible
with the following finite event order. First order the tree nodes
topologically, so that each ancestor precedes its descendants.
Place all copies for a node before that node's measurement, and place
that measurement before the copies for every subsequent node.
After all measurements, place the aligned trials in order, with their
blocks in order and all copies for a block between its neighboring
boundaries. Checkpoints and trial starts and ends are interposed in
the indicated order. The extra ordering between unrelated nodes is
harmless. In particular, every measurement copy sees its ancestor
measurements before the cut and its own and descendant measurements
after the cut; every aligned trial follows the entire tree.

\subsection*{Separated differences with prescribed residues}
We shall repeatedly use the following elementary construction.
Given a finite ordered roster, a modulus \(n\ge1\), an assigned residue
for every roster slot, and an integer \(S\ge1\), there are increasing
integer positions for the slots with those residues such that all
nonzero ordered differences are distinct and are separated from each
other and from zero by more than \(S\).  Any fixed finite number of
interposed events can also be accommodated in each prescribed gap.

Here is a direct induction.  Suppose the already chosen positions are
\(x_1<\cdots<x_a\), with all their positive differences separated by
more than \(S\), and each positive difference greater than \(S\).
Their largest positive difference is \(D=x_a-x_1\), taking \(D=0\)
when \(a=1\).  Choose \(x_{a+1}\) in its prescribed residue class so
large that
\[
 x_{a+1}-x_a>D+S,
\]
and also so that the gap has room for every event to be inserted
there.  The new positive differences all exceed \(D+S\).  The
difference of two of them is an old point gap, and is greater than
\(S\).  They are consequently separated from one another and from all
old differences.  Their negatives have the same property, and the
positive and negative differences are separated through zero.
The first gap can be chosen greater than \(S\), initiating the
induction.  Arbitrarily large representatives of every residue class
exist, so no residue restricts this choice.  Events are then placed
at distinct integer positions inside the reserved gaps.  Extra room
before the first and after the last cut accommodates any exterior
trial boundaries.

We shall choose each stencil together with its interposed events in
this way.  Denote its diameter by \(h_i\).  A subsequent translation
of the entire stage changes none of its differences or internal
chronology.  In particular, if before translation its residues are
\(v_p\), then afterwards they are \(v_p+c\) for one constant \(c\);
all differences \(v_{p'}-v_p\) used by the clock remain unchanged.

For every such stencil,
\( |\Delta_i|\le N_i^2\) and
\( |\Delta_i+\Delta_i|\le N_i^4\).
The cardinality of a sumset is at most the product of the
cardinalities of its summands, whence
\begin{equation}\label{eq:parameter-gamma-size}
 |\Gamma_i|\le\prod_{j<i}N_j^4,
 \qquad
 \Gamma_i\subseteq
 \left[-2\sum_{j<i}h_j,\,2\sum_{j<i}h_j\right].
\end{equation}
All these sets contain zero.  The second formula is a bound on the
absolute values of their elements, not a bound by the same number on
their diameters.

\subsection*{Stage 0: the number of copies and the first clock}
Put \(J=|V|^2\) and \(K=m^{J+1}\).  Initially regard \(\mathcal P_0\)
as an indexed roster, without numerical positions.  With
\(N_0=\kappa_0\mu_0\), it has \(Q=N_0(N_0-1)\) ordered pairs of
distinct slots.  For every ordered list of \(K\) distinct such pairs,
allocate a group of \(K\) main trials, whose predictions exhaust
\(M^{J+1}\).  Include one spare main trial.  Thus the number of main
trials is
\[
 T_1=1+K(Q)_K\le1+KN_0^{2K}.
\]
The full first tag alphabet consists of trial, block, and two vectors
in \(V\).  Consequently
\begin{equation}\label{eq:parameter-main-tags}
 m_1=J|V|^2T_1=J^2T_1
       \le J^2(1+KN_0^{2K}),
 \qquad |\mathcal Z|=(J+2)T_1.
\end{equation}
No position or clock modulus occurs in this count.

For completeness, the simultaneous requirement
\(\mu_0>K+c_{\rm cl}\log(2m_1)\) has the following explicit solution.
Define constants, independently of \(\mu_0\), by
\[
 A=K+c_{\rm cl}\log\bigl(2J^2(1+K\kappa_0^{2K})\bigr),
 \qquad B=2c_{\rm cl}K.
\]
Choose an integer
\begin{equation}\label{eq:parameter-copy0}
 \mu_0>\max\{1,2A,4B^2\}.
\end{equation}
For \(\mu_0\ge1\), Equation~\eqref{eq:parameter-main-tags} gives
\[
 K+c_{\rm cl}\log(2m_1)
 \le A+B\log\mu_0
 \le A+B\sqrt{\mu_0}<\mu_0.
\]
We used \(\log x\le\sqrt x\) for \(x\ge1\), which follows, for
example, by minimizing \(\sqrt x-\log x\); its minimum is
\(2-\log4>0\).  This proves the desired strict inequality without a
fixed-point assumption about the clock size.

Fix this roster, tag alphabet, prediction schedule, and
\(\mathcal Z\).  Apply Lemma~\ref{lem:clock} with this roster and
prime threshold \(1\), obtaining \(n_1,W_1\) and roster residues.
Set \(R_1=n_1\).  Reserve the prescribed residues for the stage-0
positions.  Their actual integer values will be chosen only after
the common main-trial template length is known.

\subsection*{Stage 1 and the inner marker system}
Choose
\(\mu_1=\lfloor C_1\rfloor+1\), with \(C_1\) as in
Equation~\eqref{eq:parameter-c1}; all its inputs are now fixed.  Construct
the stage-1 stencil and event order with separated differences, using
modulus \(1\), and record \(h_1\).  Choose integers
\begin{equation}\label{eq:parameter-inner-marker}
 d_\circ>10(2h_1+2R_1+5),
 \qquad K_\circ>3d_\circ^2,
\end{equation}
and obtain the inner marker rule from Lemma~\ref{lem:markers}, with
an integer inspection radius \(r_\circ\ge K_\circ\).

To prove the stage-1 count, fix \(\gamma\in\Gamma_1\) and a boundary
label \(\lambda\).  As the nonzero difference \(\delta=p-p'\)
varies, the center
\(t+\gamma+o_\lambda+\delta\) lies in an interval of diameter
\(2h_1\).  If the first-marker search changes under some shift of
absolute value at most \(R_1\), a marker lies within \(R_1\) of one
of the two original search endpoints.  For either endpoint, the
swept interval has diameter \(2h_1+2R_1<d_\circ\), so contains at
most one marker.  Such a marker accounts for at most \(2R_1+1\)
integer values of \(\delta\).  This is also the endpoint-crossing
estimate of Lemma~\ref{lem:search-crossings}; it does not depend on
the as yet unchosen halfwidth \(H\).  Distinct nonzero differences
identify distinct ordered pairs.  Summing over the two endpoints,
\(\gamma\), and \(\lambda\) gives
\[
 \#\{(p,p'):p\ne p',\ A_1^{\Gamma_1}(t+p-p',q)\}
 \le2(2R_1+1)|\mathcal Z|N_0^4<C_1.
\]
The choice of nominal boundary positions will not affect this count.

\subsection*{Stage 2 and the second clock}
The full second tag alphabet is now fixed.  Its fields have sizes
\( |\mathcal Z|\), \(|U_0|\), \(|U_0|\), \(m\), \(n_1\), and
\(|\Sigma|^{2r_\circ}\), respectively.  We take the whole Cartesian
product, including tags that no particular word realizes.  Therefore
\begin{equation}\label{eq:parameter-second-tags}
 m_2=|\mathcal Z|\,|U_0|^2m n_1|\Sigma|^{2r_\circ}.
\end{equation}
Neither the second modulus nor the cut index in a periodic window is
a field.  Choose \(\mu_2=\lfloor C_2\rfloor+1\), and apply
Lemma~\ref{lem:clock} to the resulting stage-2 roster, with threshold
\(\max\{n_1,N_M,d_\circ\}\), where \(N_M=m!\).  This gives
\(n_2,W_2\) and the stage-2 roster residues.  Construct its numerical
stencil with these residues and separated differences, and record
\(h_2\).  Set
\begin{equation}\label{eq:parameter-r2}
 R_2=n_2n_1N_Md_\circ.
\end{equation}
For a fixed \(\gamma\in\Gamma_2\), the stage-2 predicate on an
ordered pair is exactly
\[
 q-t-\gamma+p'-p\in W_2.
\]
Here \(q-t-\gamma\) is fixed while the pair varies.  The roster bound
therefore gives at most \(c_{\rm cl}\log(2m_2)\) pairs for this
\(\gamma\).  Equation~\eqref{eq:parameter-gamma-size} proves the bound
\(C_2\). The clock estimate is uniform in the translate \(q-t-\gamma\);
only \(|\Gamma_2|\), not its eventual numerical values, enters this bound.

Every prime factor of \(n_2\) exceeds \(n_1,N_M,d_\circ\).
For each \(1\le d<d_\circ\), it follows that
\[
 \gcd(n_2,dN_Mn_1)=1.
\]
Thus the cuts with displacements
\(a dN_Mn_1\), \(0\le a<n_2\), run through every second-clock
residue.  Each displacement is less than \(R_2\), as required by
Lemma~\ref{lem:periodic-split}.

\subsection*{Stage 3 and the inner search halfwidth}
Choose \(\mu_3=\lfloor C_3\rfloor+1\), construct the stage-3 stencil
with separated differences, and record \(h_3\).  Choose an integer
\begin{equation}\label{eq:parameter-h}
 H>10\bigl(h_3+R_2+r_\circ+d_\circ(1+m)+1\bigr).
\end{equation}
Fix \(\gamma\in\Gamma_3\) and \(\lambda\). The relevant centers lie in an
interval of diameter \(2h_3\). Apply Lemma~\ref{lem:search-crossings}(ii)
with \(D=2h_3\), \(L=H\), and \(h=2R_2\).
Equation~\eqref{eq:parameter-h} gives \(D<2L+2\).
At most \(2R_2\) centers have the required empty search and right-margin
marker. Distinct ordered differences identify distinct ordered pairs,
and a union bound gives
\[
 \#\{(p,p'):p\ne p',\ A_3^{\Gamma_3}(t+p-p',q)\}
 \le2R_2|\mathcal Z|\prod_{j<3}N_j^4<C_3.
\]
This argument is independent of the nominal positions of the
boundaries.  Equation~\eqref{eq:parameter-h} also implies the separate
periodic-window requirement
\begin{equation}\label{eq:parameter-periodic-margin}
 H-d_\circ>2r_\circ+m d_\circ+R_2.
\end{equation}
Indeed, subtracting the right side from the left and using
Equation~\eqref{eq:parameter-h} leaves a positive quantity.  Together
with \(K_\circ>3d_\circ^2\), this supplies both the overlapping
short-period blocks and all the context and stabilized-power
margins used in Lemma~\ref{lem:periodic-split}.

At this point the first three numerical stencils, both clocks, the inner
marker radius, and the periodic-window width have all been fixed. The
only postponed early choice is the numerical stage-0 stencil. We can
now place it with gaps exceeding the complete main-trial template while
preserving the first-clock residues chosen earlier.

\subsection*{The main template and the postponed stage-0 positions}
Choose an even integer
\[
 G>10(H+r_\circ+R_1+R_2+1),
 \qquad T=(J+3)G.
\]
The common template is the interval of length \(T\).  Boundary
\(i\), \(0\le i\le J+1\), has center \((i+1)G\) relative to its
start; block \(j\), \(1\le j\le J\), begins at
\(jG+G/2\) and has the \(R_1\) consecutive cut positions prescribed
in Section~\ref{sec:outer}.  All boundary inspection windows, and
all periodic-window cuts and their inspection neighborhoods, fit
strictly inside the template.  For a block cut, its preceding
boundary's entire inspection window lies to the left, and the next
boundary's entire inspection window lies to the right.
The same holds for the future searches under every origin
\(k-c_{\nu,j}-a\), \(0\le a<R_1\): relative to \(k\) the nearest
future center is \(G/2-a\), which exceeds \(H+r_\circ\).
The nearest preceding center for the actual prefix is
\(-G/2-a\), whose whole inspection window precedes the cut.
Earlier and later trial windows have still larger separations.
The periodic suffix's absence test uses starts from
\(k-r_\circ\) onward, so a marker inspection needs no letter earlier
than \(k-2r_\circ\).  This is exactly the context length included
in its tag.

We may now choose the numerical positions of the stage-0 roster,
with its previously prescribed first-clock residues.  Use the
separated-difference induction with spacing \(S=T\), reserving all
aligned events in their required order.  Record \(h_0\).
For each allocated group, let
\(\delta_a=p_a-p'_a\), \(1\le a\le K\), be the differences of its
listed pairs.  Give its trial \(a\) the template start
\begin{equation}\label{eq:parameter-translated-template}
 b_{\rm group}-\delta_a
\end{equation}
relative to \(t\), retaining the prediction already assigned to
that abstract trial.  The differences are distinct and separated by
more than \(T\), so the templates within a group are disjoint.
Bases for distinct groups can be chosen explicitly.  Put the spare
trial at offset \(T\), enumerate the groups by \(g=1,\ldots,(Q)_K\),
and take
\[
 b_g=T+g(2h_0+2T)+h_0.
\]
Every group lies between offsets
\(T+g(2h_0+2T)\) and
\(T+g(2h_0+2T)+2h_0+T\).  These intervals are positive, mutually
separated, and disjoint from the spare template.  Sort all trials by
physical position for the processing rules.  Sorting changes neither
their identities nor their assigned predictions.  It now determines
all \(o_\lambda,c_{\nu,j},\mathcal H_1,\mathcal H_2\).

It remains to justify \(C_0\).  Its first-clock part contributes at
most \(c_{\rm cl}\log(2m_1)\), since
\[
 q-(t+p-p')=q-t+p'-p
\]
and the first-clock roster differences have been preserved.
Suppose that the prediction-failure part held for \(K\) distinct
ordered pairs.  Take their ordered list and its allocated group.
At origin \(t+\delta_a\), the indicated trial has physical template
start
\[
 (t+\delta_a)+(b_{\rm group}-\delta_a)=t+b_{\rm group}.
\]
All these \(K\) trials therefore have identical physical boundary
search centers.  The predicate under consideration includes
\(\mathcal A(t+\delta_a)\), so all their boundaries exist.
Locality and translation equivariance of the marker rule give the
same actual boundary positions, hence the same actual tuple of
\(J+1\) interval products, for all \(K\) trials.  Their predictions
exhaust all possible tuples; one must be perfect.  This contradicts
the assumed absence of any perfectly predicted main trial at its
origin.  There are thus at most \(K-1\) prediction-failure pairs,
and the total is strictly less than \(C_0\).  If the roster has fewer
than \(K\) ordered pairs, that last assertion already follows from
its size.  The internal visibility needed for this argument, including
all translated templates, is established below.

\subsection*{Stage 4 and the end marker system}
Choose an integer
\begin{equation}\label{eq:parameter-we}
 w_E>1+\max\left\{
 \operatorname{diam}\mathcal H_1,
 \operatorname{diam}\mathcal H_2,
 2(h_0+h_1+h_2+h_3)\right\}.
\end{equation}
In particular \(h_i<w_E\) for \(i<4\).  Choose
\(\mu_4=\lfloor C_4\rfloor+1\), construct its stencil with separated
differences, and record \(h_4\).  Choose integers
\begin{equation}\label{eq:parameter-end-marker}
 d_\bullet>10(2h_4+2w_E+5),
 \qquad K_\bullet>3d_\bullet^2,
\end{equation}
and obtain its marker rule and radius \(r_\bullet\ge K_\bullet\).
These choices imply
\[
 w_E<d_\bullet,
 \qquad K_\bullet-w_E>d_\bullet^2,
\]
so Lemma~\ref{lem:shared-anchor} applies to each required set of
unexpanded candidate origins.
For fixed \(\gamma\in\Gamma_4\), apply the same endpoint argument as
at stage 1, with center range of diameter \(2h_4\), stability radius
\(w_E\), and marker separation \(d_\bullet\).  Each endpoint sweep
has diameter \(2h_4+2w_E<d_\bullet\), and the number of bad
nonzero differences is at most \(2(2w_E+1)\).  Hence
\[
 \#\{(p,p'):p\ne p',\ A_4^{\Gamma_4}(t+p-p',q)\}
 \le2(2w_E+1)\prod_{j<4}N_j^4<C_4.
\]
This count is independent of the end search's eventual lag.

All five ambiguity bounds are now proved, and each stencil has more
copies per kind than its bound. It remains to put the finite arrangements
inside words. This last step must cover every false-origin inspection,
not only the windows used by a correctly chosen split.

\subsection*{Final placement and visibility under every hypothesis}
Translate the five complete aligned arrangements, including all
interposed events, into disjoint integer intervals strictly after
\(s\), in order \(4,3,2,1,0\).  End stage \(i\)'s interval at
\(a_i\), after all its events and before the next interval begins.
Write \(A_{\rm buf}\) for a positive upper bound on all their offsets
from \(s\).  These translations do not change any \(h_i,N_i\),
\(\Delta_i,\Gamma_i\), or clock difference.  Every stage-\(i\)
measurement, trial, and endpoint is after \(a_{i+1}\), for \(i<4\).
This is the chronological hypothesis for the recursion of the total
suffix tables in Lemma~\ref{lem:total-tables}.

We give explicit finite envelopes, so that the remaining choices do
not hide a dependence on an episode length.  Let
\begin{equation}\label{eq:parameter-origin-envelope}
 D=1+\max\left\{
 \operatorname{diam}\mathcal H_1,
 \operatorname{diam}\mathcal H_2,
 2R_1,\ 2\sum_{i=0}^4h_i\right\}.
\end{equation}
Every origin displacement that must be tested, before the extra
stability shifts in \(A_1\) and \(A_4\), has absolute value less
than \(D\).  To check this assertion exhaustively, an outer end guard
uses differences in \(\mathcal H_j-\mathcal H_j\).  At an actual
main cut, even a suffix with an unrelated trial/block role uses a
candidate \(k-c_{\nu,j}-a\) whose displacement from \(t\) is in
\(\mathcal H_1-\mathcal H_1\).  Aligned candidates have displacements
in \(\Delta_i\).  Current-witness and primary later-witness tests
use \(\Delta_i+\Gamma_i\); the secondary later-witness tests use
\(\Delta_i+\Gamma_i+\Delta_i\).  Designations and the next-domain
promises use subsets of these same expansions, or
\(\Gamma_{i+1}=\Gamma_i+\Delta_i+\Delta_i\).
Equation~\eqref{eq:parameter-gamma-size} bounds all their absolute
values by \(2\sum_{j\le i}h_j\).  The swept intervals used in the
count proofs satisfy the same bounds even when the displacement
ranges over every integer in \([-h_i,h_i]\).
The additional shifts of \(A_1\) have magnitude at most \(R_1\),
and those of \(A_4\) at most \(w_E\); we include them in the respective
inspection envelopes below.

Let \(A_{\rm main}\ge1\) be an integer exceeding the absolute values
of all main-template endpoints, all \(o_\lambda\), and all elements
of \(\mathcal H_1\cup\mathcal H_2\).  Set
\begin{equation}\label{eq:parameter-inner-envelope}
 B_\circ=A_{\rm main}+D+H+2R_2+R_1+r_\circ+K_\circ+2.
\end{equation}
All bounded inner-marker, boundary, main-trial, and periodic-window
data at every origin just listed lie strictly between
\(t-B_\circ\) and \(t+B_\circ\).  This includes inspection
neighborhoods of radius \(r_\circ\), the right-margin overhang
\(2R_2\), the preceding periodic contexts, the short-period blocks,
and the whole swept ranges used in the estimates.
Equation~\eqref{eq:parameter-h} gives \(H>r_\circ\), so the term
\(H+r_\circ\) also covers the \(2r_\circ\)-letter context.
Actual marker-selected interval endpoints differ from their centers by at
most \(H\), so products within these bounded trials have endpoints
in the same envelope.  Choose
\begin{equation}\label{eq:parameter-l0}
 L_0>A_{\rm buf}+B_\circ+2.
\end{equation}
Consequently this entire envelope follows every aligned buffer.  It
is available in every aligned suffix as far as its left endpoint is
concerned, including before time has been decoded.  Its right
endpoint will be put before every possible true anchor next.

For the end controls put
\begin{equation}\label{eq:parameter-end-envelope}
 B_\bullet=D+w_E+d_\bullet+r_\bullet+K_\bullet+2.
\end{equation}
Every bounded end-marker inspection and seed block required under
any candidate origin or witness translate lies strictly between
\(z_\bullet(t)-B_\bullet\) and
\(z_\bullet(t)+B_\bullet\).
Indeed the origin contributes at most \(D\), a stability test at most
\(w_E\), a searched marker start at most \(d_\bullet\), and its
inspection at most \(r_\bullet\); a seed extends by at most
\(K_\bullet\).  We have used a sum of these bounds to cover all
cases at once.  Choose
\begin{equation}\label{eq:parameter-final-margins}
 o_\bullet>B_\circ+B_\bullet+2,
 \qquad B_{\rm tail}>B_\bullet+d_\bullet+2.
\end{equation}
The entire bounded end envelope now follows the inner envelope and
all actual cuts.  Even an end-rule origin among the stated
hypotheses has its marker anchor, or the beginning of its possible
first-break scan, after the earlier activities, since
\(B_\bullet>D+d_\bullet\).
For the true episode, its anchor always satisfies
\(q\ge z_\bullet(t)-d_\bullet\).  Therefore
\[
 q+B_{\rm tail}>z_\bullet(t)+B_\bullet+2.
\]
Every bounded control is thus strictly internal to the true episode,
even when its anchor is the earliest possible marker.  In the
markerless case the actual first break may occur arbitrarily later;
this only increases the available length.  Its mismatching letter
is visible because \(B_{\rm tail}>1\).  In the notation of Equation~\eqref{eq:geometry-margins}, take
\(I_-=-B_\circ,I_+=B_\circ,E_-=-B_\bullet,E_+=B_\bullet\)
and use \(A_{\rm buf}\) for its last buffer endpoint.  Equations~\ref{eq:parameter-l0}
and~\ref{eq:parameter-final-margins} then imply each of those
margin inequalities directly.

In particular, the stage-4 controls do not presuppose knowing the
end.  At \(k=s+p\), all their inspected letters have displacement
from \(z_\bullet(t)\) at most
\[
 h_4+2\sum_{j<4}h_j+w_E+d_\bullet+r_\bullet+1.
\]
A sufficient padding inequality just for these controls is
\begin{equation}\label{eq:parameter-stage4-padding}
 B_{\rm tail}>
 h_4+2\sum_{j<4}h_j+w_E+2d_\bullet+r_\bullet+1,
\end{equation}
which follows from Equations~\ref{eq:parameter-origin-envelope},
\ref{eq:parameter-end-envelope}, and~\ref{eq:parameter-final-margins}.
It contains no occurrence of the unbounded break position \(q\).

We spell out the consequences for arbitrary prefix/suffix pairings.
Every proposed suffix must have all of its specified bounded
controls on its reading side; an unavailable candidate is a reason
to reject the suffix, not a reason to remove that candidate.
At a cut of a true complete episode, a shorter proposed suffix either
fails this availability test or reads exactly the true letters in
all those windows.  A longer one reads the same letters there,
because all the windows precede the true first episode's end.
Locality of both marker rules therefore fixes the same control values
in either case.  This applies to every candidate, including the true
candidate, independently of the proposed suffix length.
For \(i<4\), the candidate-origin diameter is \(h_i<w_E\); for the
outer guards it is at most the diameter of the appropriate
\(\mathcal H_j\).  The equal-search and seed-overlap argument of
Lemma~\ref{lem:shared-anchor} hence supplies the true anchor before
any anchor-dependent witness is used.  At stage 4 the witness is
\(q\)-free: the true candidate cannot disappear from the uniqueness
test, because its entire inspected neighborhood is already internal
by Equation~\eqref{eq:parameter-stage4-padding}.  Unique decoding
therefore precedes the exact end test at that decoded origin, just
as required by Lemma~\ref{lem:time-decoder}.

Finally, the side of every noncontrol computation is determined by
the event order.  Prefixes read earlier measurements, earlier trial
boundaries, and the indicated prefix products.  Suffixes at aligned
trial cuts read the later boundaries, the later trials, and the
terminal table at \(a_i\).  At a measurement copy they read its
measurement and every descendant measurement after the cut, using
ancestor values only as verified metadata.  The reverse stage order
puts every earlier-stage event needed by a terminal table after the
current buffer.  The main-template inequalities already put all
main prefix and suffix boundary searches on their specified sides,
including all suffix role hypotheses.  The periodic suffix uses
only its stated \(2r_\circ\) letters of supplied left context.
After correct decoding, all required suffix computations thus have
their prescribed access.  For an arbitrary standalone suffix with
false metadata, a needed unavailable position or ill-ordered
interval is rejected; its evaluation never calls for an unspecified
prefix factor.  Products that extend to the true or proposed end
have that end as an endpoint, and the total formulas \(\Phi_i\)
use the supplied hypothetical product for their omitted prefix.

These observations also give the finite-endpoint property used in
Section~\ref{sec:compilation}.  In any single test, all local
inspections and all interval endpoints other than its exact end
range over fixed finite sets of offsets from the beginning of that
test; bounded offsets from its exact end are allowed as well.
There is only the one first-break scan belonging to its selected
reference or decoded end-rule origin.  Witness translations and
hypothetical table evaluations keep that scan's \(q\) fixed and
introduce no additional search for a break.  This assertion holds
for false suffix metadata as well as for successful pairings.

Later translations preserve every earlier count and residue
difference. The final three lengths change no schedule diameter,
clock, state set, or number of copies. Every marker count now applies
to complete words because its full inspection windows are internal,
including those at false origins. Thus the ambiguity bounds,
chronological requirements, and end margins in the proposition all hold.
\end{proof}
